\section{Introduction}

Forensic analytics prioritize investigation. They do not prove fraud, certify
outcomes, or replace risk-limiting audits. A batch/scanner effect is a
systematic deviation associated with a tabulation batch, scanner, drop box,
upload, or similar operational grouping. It is a grouping-level signal: the
question is whether units processed through the same group move together after
ordinary comparisons are accounted for, not whether any single ballot is
suspicious.

This boundary also separates the W-series from the V-series. V-series work
replays audit or canvass evidence and can support certification-style claims
when its assumptions are met. W-series work produces investigative analytics:
outlier scores, residuals, digit checks, batch effects, spatial anomalies, and
discrepancy clustering. The output is an investigation queue with caveats, never
a pass/fail certification result.

Batch and scanner effects are a natural W-series family because election data
often arrive in operational groups. Vote-by-mail batches may differ from
election-day batches. Early-vote uploads may differ from precinct uploads.
Scanners may serve different locations, times, or ballot styles. Some group
differences are therefore real and innocent. The method in this paper is meant
to identify groups whose deviations deserve review while preserving that
interpretive boundary.
